\chapter{Ameaças à Validade}
\label{chap:ameacas_validade}

Toda pesquisa empírica possui restrições e potenciais fontes de viés. A seção de Ameaças à Validade demonstra a maturidade científica do autor ao reconhecer de forma transparente os fatores que podem limitar a abrangência ou a confiabilidade das conclusões.

Na engenharia de software e ciência da computação, adota-se a classificação proposta por Wohlin et al., dividida em quatro dimensões fundamentais:

\begin{itemize}
    \item \textbf{Validade Interna:} Trata de variáveis de confusão ou fatores não controlados que poderiam ter influenciado os resultados observados (ex.: variação de hardware, viés de seleção de amostras). Deve-se descrever as estratégias de mitigação aplicadas.
    \item \textbf{Validade Externa:} Analisa o grau de generalização dos achados para outros contextos, linguagens de programação, magnitudes de dados ou domínios de aplicação.
    \item \textbf{Validade de Construção:} Examina se as métricas e ferramentas utilizadas de fato medem os conceitos teóricos que se pretendia avaliar.
    \item \textbf{Validade de Conclusão Estatística:} Avalia a adequação dos testes estatísticos empregados, o tamanho da amostra e o poder de significância dos dados coletados.
\end{itemize}

O mapeamento honesto dessas ameaças fortalece a credibilidade da pesquisa e fornece subsídios para o planejamento dos Trabalhos Futuros.